\chapter{Conclusiones y Trabajo Futuro}
\label{cap:conclusiones}

Este capítulo cierra la memoria. Primero repasa los objetivos de la introducción uno a uno, con la evidencia que sostiene cada uno y lo que quedó a medias; después recoge las limitaciones del sistema tal como se entrega; y por último agrupa el trabajo futuro por capas, con la razón por la que cada línea no entró en el TFG.

\section{Conclusiones}
\label{sec:conclusionesObjetivos}

El objetivo general era diseñar, construir y evaluar una plataforma web para nutricionistas cuyo núcleo fuese un copiloto de generación de planes. La plataforma existe y está desplegada (sección~\ref{sec:despliegue}): el profesional escribe lo que necesita en castellano, el traductor lo convierte en restricciones formales, el motor construye un borrador que las cumple, el validador lo audita y el plan solo llega al cliente cuando el profesional lo firma (capítulos~\ref{cap:disenoSistema} a~\ref{cap:copilotoIA}). El proceso completo de producto también se recorrió: la idea se refinó con nutricionistas antes de escribir código, el alcance se reformuló con lo que pedían, y el resultado se midió con criterios fijados antes de medir (capítulo~\ref{cap:evaluacion}). La pieza que falta de ese recorrido es la prueba con nutricionistas sobre la aplicación desplegada, que está en curso en el momento de redactar esta versión (sección~\ref{sec:evalUsuarios}).

Respecto a los objetivos específicos:

\begin{enumerate}
    \item \textit{Analizar el flujo de trabajo real de un nutricionista para extraer los requisitos.} Los requisitos salieron de las conversaciones con el nutricionista que motivó el proyecto y con los profesionales que contacté después (sección~\ref{sec:reformulacionAlcance}), y el catálogo inicial de 61 requisitos se redujo a 27 cuando esas mismas conversaciones enseñaron dónde estaba el cuello de botella. Ese recorte es la parte del objetivo que más valor tuvo: el copiloto no estaba en la idea original y salió de escuchar la demanda real. Dos requisitos conservados en el capítulo~\ref{cap:analisisRequisitos} no se llegaron a construir (las notas privadas del cliente, RF-03, y el \textit{dashboard} de métricas del copiloto, RF-39); se recogen en la sección~\ref{sec:limitaciones}.

    \item \textit{Realizar un estudio de mercado y análisis de la competencia.} El capítulo~\ref{cap:estadoDeLaCuestion} analiza cuatro plataformas del sector con sus webs consultadas en septiembre de 2026 y las compara con lo que la aplicación entrega (tabla~\ref{tab:competidores}); la comparativa de la primera versión, más amplia y anterior a la reformulación, se conserva en el Anexo~\ref{apendice:catalogoV1}. El posicionamiento que resulta es concreto: las plataformas que generan planes lo hacen a partir de números y ninguna explica por qué un plan no puede existir (sección~\ref{sec:posicionamiento}). El análisis se quedó a propósito en el posicionamiento funcional; un estudio económico del producto queda fuera de este trabajo.

    \item \textit{Diseñar la arquitectura del sistema y su modelo de datos.} Los dos paneles comparten una misma base de datos y el aislamiento entre profesionales y entre cliente y profesional está en la propia capa de datos, mediante políticas de seguridad a nivel de fila que los runners multi-actor comprueban tabla a tabla (secciones~\ref{sec:capaDatos}, \ref{sec:capaSeguridad} y~\ref{sec:evalBancos}). La tabla de restricciones funciona como interfaz entre tres productores (el formulario del cliente, el traductor y el estilo clínico del profesional) y un único consumidor, el motor (figura~\ref{fig:productorConsumidor}); esa separación es la que permitió construir y verificar el motor antes de que existiera la capa de IA. Este objetivo se cumplió tal como estaba planteado.

    \item \textit{Construir un motor de generación basado en programación con restricciones, con un validador independiente y sin depender de la IA.} El motor del capítulo~\ref{cap:motorGeneracion} genera, valida, persiste y firma planes sin ningún modelo de lenguaje, y así se verificó primero con los bancos deterministas y después de extremo a extremo contra el despliegue (sección~\ref{sec:evalBancos}). La calidad de los planes se midió y se corrigió por capas: la auditoría de sinsentidos bajó las raciones en los límites del 66\,\% al 6\,\% (sección~\ref{sec:evalAuditoria}), y la función objetivo se reequilibró dos veces con datos (secciones~\ref{sec:evalMagnitudes} y~\ref{sec:evalBandas}). Quedaron límites conocidos: las reglas estructurales se asumen válidas para la población general y el profesional no puede anularlas (sección~\ref{sec:estructurales}), el término de reparto por comidas sigue dominando la función objetivo cuando está activo (sección~\ref{sec:objetivoResolucion}) y con el catálogo de cien alimentos el solver no siempre demuestra el óptimo dentro del límite de tiempo (sección~\ref{sec:evalEscalado}).

    \item \textit{Integrar la capa de IA que hace usable el motor en lenguaje natural.} El traductor, la revisión previa del mensaje, la generación asíncrona y la pantalla de arbitraje están construidos e integrados (capítulo~\ref{cap:copilotoIA}), y el profesional decide sobre cada restricción propuesta y sobre cada plan antes de que llegue al cliente. El modelo se ejecuta en local para que el dato clínico no salga a ningún proveedor, y la evaluación enseñó lo que eso cuesta: un modelo de siete mil millones de parámetros comete fallos residuales que el diseño contiene pero no elimina (sección~\ref{sec:evalModelosLocales}), y el estudio de formatos descartó con datos construir un módulo que reescribiera el mensaje (sección~\ref{sec:evalFormatos}). La parte que quedó a medias es el estilo clínico: se declara como reglas, y aprenderlo de las correcciones del profesional no se abordó.
\end{enumerate}

La prueba con nutricionistas (sección~\ref{sec:evalUsuarios}) cerró el ciclo con dos profesionales: las ocho tareas les resultaron fáciles, destacaron lo intuitivo de la aplicación y el tiempo que ahorra, y su único pero fue la espera al generar la dieta. Sus dos pautas, un worker siempre disponible y una prueba con más profesionales, están en la sección~\ref{sec:trabajoFuturo}.

Hay tres cosas que aprendí en el proceso y que condicionaron el diseño más que cualquier tecnología concreta. Cerrar el vocabulario de restricciones antes de tocar el modelo de lenguaje hizo que la IA fuese una capa de conveniencia sobre un motor que ya funcionaba, y no al revés. Verificar cada capa por separado, con bancos que se re-ejecutan al cerrar cada bloque, es lo que permitió cambiar el motor varias veces en verano sin romper lo anterior. Y fijar el criterio antes de medir dio el mismo valor a los resultados negativos que a los positivos: el reescritor y la rotura de simetría no se construyeron porque los datos no los respaldaban, y las bandas de tolerancia se mantuvieron con la desviación del criterio documentada (sección~\ref{sec:evalBandas}).

\section{Limitaciones}
\label{sec:limitaciones}

Las limitaciones que siguen son las del sistema tal como se entrega. Casi todas están ya dichas en el capítulo donde se tomó la decisión; aquí se juntan para que el lector las tenga en un solo sitio.

En el motor, las reglas estructurales se asumen válidas para la población general y el profesional no puede anularlas caso a caso, así que una prescripción fuera de esos rangos, como la dieta renal del ejemplo de la sección~\ref{sec:estructurales}, se declararía infactible. El validador solo rechaza lo que el solver garantiza y avisa del resto, con lo que un plan con un micronutriente por encima del tope llega al profesional como aviso y no como rechazo (sección~\ref{sec:validador}). El término de reparto por comidas domina la función objetivo cuando está activo, el cociente entre nutrientes solo existe como restricción dura (sección~\ref{sec:objetivoResolucion}) y una restricción dura de energía o de macronutriente exige la igualdad exacta por día, sin la tolerancia que sí tienen las blandas. Con el catálogo de cien alimentos el solver no siempre demuestra el óptimo dentro del límite de tiempo, y ese límite es de 90 segundos con dos hilos porque el alojamiento gratuito no da para más (secciones~\ref{sec:evalEscalado} y~\ref{sec:evalHosting}).

El catálogo es propio: cien alimentos con sus nutrientes y perfiles de ración cargados a mano, sin la base de referencia BEDCA que preveía el capítulo~\ref{cap:analisisRequisitos}. Los alimentos que añade el profesional no tienen perfil de ración y el motor les aplica un rango genérico (sección~\ref{sec:plausibilidad}); y los umbrales de los casos de infactibilidad de los bancos se han tenido que recalibrar cada vez que el catálogo creció, porque con más fuentes de un nutriente cuesta más demostrar que algo no puede existir.

En el copiloto, el modelo de lenguaje corre en local y es pequeño; la evaluación enseñó fallos residuales que el diseño contiene pero no elimina, como confundir una categoría con un alimento concreto (sección~\ref{sec:evalModelosLocales}). El worker que ejecuta ese modelo es, durante el desarrollo y la evaluación, mi propio ordenador: si está apagado, la aplicación funciona entera y las tareas de IA esperan en la cola (sección~\ref{sec:despliegue}).

En la plataforma, la interfaz está en inglés aunque la petición al copiloto se escriba en castellano, porque los mockups se diseñaron así desde el principio. Dos requisitos que el capítulo~\ref{cap:analisisRequisitos} conserva no se construyeron: las notas privadas del profesional sobre cada cliente (RF-03) y el \textit{dashboard} de métricas del copiloto (RF-39), que en su lugar enseña datos reales de clientes, planes y alimentos. Algunas pantallas tienen recortes frente al diseño, la mensajería se refresca por sondeo cada quince segundos, las altas que escriben en varias tablas desde el navegador no son transaccionales y se compensan con un borrado lógico si falla un paso (sección~\ref{sec:capaDatos}), los datos de demostración llevan fechas relativas que caducan (sección~\ref{sec:validacionTesting}) y el repositorio no tiene protección de rama.

En la evaluación, la prueba con nutricionistas se hizo con dos profesionales y sin comentarios por tarea, así que valida los flujos pero no mide la usabilidad con detalle (sección~\ref{sec:evalUsuarios}); los bancos del motor usan casos construidos a mano, no planes reales de profesionales; y la comparativa con los competidores se apoya en lo que sus webs decían en septiembre de 2026, que cambiará.

\section{Trabajo futuro}
\label{sec:trabajoFuturo}

Las líneas se agrupan por la capa a la que pertenecen. Cada una lleva la razón por la que no entró en el TFG: tiempo, una decisión de alcance tomada a propósito, o una dependencia de otra pieza que tampoco existe.

\subsection{Motor de generación}
\label{sec:futuroMotor}

\begin{itemize}
    \item \textbf{Variedad por familias de alimentos.} El término de alimentos sin usar es lineal: pasar de cinco a seis alimentos distintos vale lo mismo que pasar de cincuenta a cincuenta y uno, y un plan con ochenta alimentos a la semana no es mejor ni más fácil de seguir. La versión que tiene sentido clínico premia que cada familia esté presente con un peso propio, sin fijar un número de ingredientes. Requiere una taxonomía de familias en el catálogo que hoy solo existe en parte; mientras, la variedad la controla el estilo clínico del profesional (sección~\ref{sec:evalBandas}).
    \item \textbf{Empates entre soluciones.} Cuando varios planes empatan en la función objetivo, el solver devuelve el primero que encuentra y nada distingue a los demás. Uno de los directores del TFG señaló dos salidas: un orden de prioridad entre las restricciones, resolviendo por niveles para que el empate lo deshaga la más importante (un objetivo lexicográfico), o devolver una lista corta de planes alternativos con el mismo valor para que el profesional elija, que encaja con el marco del proyecto: el sistema propone y el profesional decide. Las dos se apoyan en el modelo actual; la primera cambia la estrategia de resolución y la segunda la interfaz. La variedad por familias del punto anterior sería un criterio natural para deshacer esos empates.
    \item \textbf{Reparto calórico por comidas.} Renormalizar el término para quitarle el factor de mil que arrastra su linealización, con división entera y una banda porcentual como la de la energía. No se midió porque solo dos casos de los bancos lo activan y uno de ellos es un conflicto deliberado (sección~\ref{sec:evalMagnitudes}).
    \item \textbf{Tolerancias configurables.} Que cada fila de objetivo calórico o de macronutriente lleve su propia banda en vez de la global, y que las filas duras tengan una tolerancia estrecha en lugar de la igualdad exacta por día: un ``debe'' clínico significa alrededor de un valor, no ese valor al gramo. Toca la validación, el formulario y el traductor a la vez; decidí no tocar el motor en el cierre del proyecto.
    \item \textbf{Reglas estructurales editables.} Que el profesional pueda ver las ocho reglas de sentido común, retirarlas o endurecerlas para un cliente concreto, y así cubrir casos como el renal sin declararlos infactibles. No hubo tiempo de hacerlas configurables; se documentaron como asunción (sección~\ref{sec:estructurales}).
    \item \textbf{Acompañantes.} Un rol de acompañante (tomate, cebolla, lechuga) que no consuma el tope de cuatro alimentos por comida, como señaló uno de los directores del TFG al revisar el tope; exige marcar ese rol en el catálogo.
    \item \textbf{Cociente entre nutrientes blando.} Hoy solo existe como restricción dura porque penalizar la desviación de un cociente obliga a linealizarlo (sección~\ref{sec:catalogoClinico}).
    \item \textbf{Recetas y sustituciones (RF-26 y RF-27).} El motor trabaja con alimentos sueltos por diseño. Las recetas exigirían que el motor eligiera entre preparaciones con su composición agregada y una interfaz para definirlas; las sustituciones, un modelo de equivalencia que ninguna pieza tiene hoy.
    \item \textbf{Validación por patología y medicación.} El validador aplica reglas de población general. Afinarlo por patología o por fármaco depende de ampliar la ficha del cliente con esas variables y el catálogo con los nutrientes implicados, ampliación que el capítulo~\ref{cap:analisisRequisitos} pospone a propósito.
    \item \textbf{Residuos de la capa de sentido común.} Un desayuno de solo frutos secos sigue siendo posible en teoría, y los tentempiés de una pieza y los alimentos sin perfil dependen de rangos genéricos. La prueba con nutricionistas no señaló ninguno de estos casos; la calibración fina queda pendiente de una prueba con más profesionales.
\end{itemize}

\subsection{Copiloto de IA}
\label{sec:futuroCopiloto}

\begin{itemize}
    \item \textbf{Aprender el estilo del profesional.} El estilo clínico se declara como reglas (sección~\ref{sec:arbitraje}). La versión ambiciosa lo aprende de las correcciones que el profesional hace sobre cada borrador y lo arranca con un cuestionario. Se construyó la versión mínima acordada con los tutores; el aprendizaje necesita además la edición del plan, que tampoco existe.
    \item \textbf{Edición del plan generado.} Hoy el borrador se acepta o se regenera con otras restricciones; falta poder cambiar un alimento o una cantidad y que el validador vuelva a auditar el resultado. Se dejó fuera porque el flujo de generación sustituyó a la edición manual de planes en el alcance.
    \item \textbf{Métricas del copiloto (RF-39).} La tasa de aprobación al primer intento y la distancia entre el borrador y la versión entregada, tal como las define el capítulo~\ref{cap:analisisRequisitos}, solo se pueden medir cuando exista esa edición; los tokens y la latencia de cada llamada ya quedan en la traza (sección~\ref{sec:eleccionLLM}).
    \item \textbf{Motor de lenguaje remoto.} El traductor habla con el modelo a través de una interfaz mínima, así que un motor remoto con garantías contractuales de no retención sería otra implementación de esa interfaz. Permitiría comparar coste, latencia y calidad frente al modelo local y repetir el estudio de formatos con un modelo mayor; no entró por la decisión de privacidad de la sección~\ref{sec:eleccionLLM} y porque el modelo local bastó para el alcance.
    \item \textbf{Reescritor del mensaje.} No es trabajo futuro: el estudio de formatos descartó construirlo con datos (sección~\ref{sec:evalFormatos}), y solo se reabriría si un modelo distinto cambiara ese resultado.
\end{itemize}

\subsection{Plataforma}
\label{sec:futuroPlataforma}

\begin{itemize}
    \item \textbf{Catálogo de referencia.} Cargar BEDCA o USDA con sus nutrientes, manteniendo los perfiles de ración por alimento que la capa de sentido común necesita. El catálogo propio de cien alimentos bastó para construir y evaluar el motor.
    \item \textbf{Funciones de gestión recortadas por la versión mínima.} Editar y borrar alimentos propios, las notas privadas del cliente (RF-03), el calendario en rejilla semanal, las notificaciones y la invitación de clientes desde el panel del profesional. Son pantallas con su tabla detrás ya existente o trivial, dejadas fuera por tiempo.
    \item \textbf{Seguimiento.} Volcar el peso que registra el cliente a la ficha con validación del profesional, y caducar las solicitudes de cita pendientes con un proceso periódico. Quedaron fuera del alcance del panel del cliente.
    \item \textbf{Mensajería en tiempo real.} El sondeo cada quince segundos fue una decisión deliberada para el volumen de un profesional; las suscripciones en tiempo real de la propia base de datos quitarían la latencia a cambio de gestionar conexiones persistentes.
    \item \textbf{Interfaz en castellano.} Traducir la interfaz, en inglés desde los mockups, es la mejora que más notarían los profesionales a los que va dirigida.
    \item \textbf{Despliegue.} Mover el worker a un servidor de la clínica o a la nube, y pasar a un alojamiento de pago que permita bajar el límite de tiempo del solver y devolverle los hilos (sección~\ref{sec:evalHosting}).
    \item \textbf{Robustez de datos.} Pasar las altas de varias tablas desde el navegador a funciones de la base de datos, para que sean transaccionales; y datos de demostración con fechas que se vuelvan a anclar solas.
\end{itemize}

\subsection{Evaluación}
\label{sec:futuroEvaluacion}

La prueba con nutricionistas (sección~\ref{sec:evalUsuarios}) deja dos pautas: que el worker de generación esté siempre disponible, porque la espera fue su único pero, y repetirla con más profesionales. Después, la evaluación que tiene más valor es la de uso real: las métricas del copiloto que define RF-39, medidas sobre planes de profesionales en vez de sobre casos construidos a mano, y la comparativa con un motor de lenguaje remoto. Y una lección de método que conviene arrastrar a cualquier experimento futuro sobre el motor: los criterios sobre conteos deben llevar una tolerancia del tamaño de la brecha de parada del solver, o exigir óptimos a brecha cero (sección~\ref{sec:evalBandas}).

\paragraph{Sobre los requisitos de la primera versión.} Los cinco bloques que el capítulo~\ref{cap:analisisRequisitos} elimina (formularios, alertas, lista de la compra, marketplace y automatizaciones) y la mayoría de los requisitos individuales de la tabla~\ref{tab:rfIndividualesFuera} no aparecen en las listas anteriores a propósito. Se descartaron porque los nutricionistas no los pedían y porque no contribuyen al copiloto, y el Anexo~\ref{apendice:catalogoV1} los conserva como registro de la idea inicial. Solo recojo como trabajo futuro los que tienen un enganche en lo construido: las recetas y sustituciones, las notas privadas y el \textit{dashboard} de métricas.

\section{Cierre}
\label{sec:cierre}

% Cierre en la voz de Victor: borrador del chat para que lo reescriba a su gusto.
Empecé este trabajo con la condición de que sirviera para algo real, y lo que más me ha enseñado es lo lejos que estaba mi primera idea de lo que los nutricionistas necesitaban. Sin las conversaciones con ellos habría construido una plataforma de gestión más, y el copiloto, la parte que da sentido al proyecto, no habría existido. La otra lección es la indicación de mis directores del trabajo de fin de grado: la versión mínima que aplique las funcionalidades, verificada, vale más que una ambiciosa a medias. Aplicarla me costó al principio, porque suponía dejar fuera cosas que ya tenía diseñadas, pero es lo que permitió llegar a una aplicación desplegada y con un motor medido en vez de a un prototipo bonito. Lo que digan los profesionales cuando la usen es lo que queda por escuchar.
